% ============================================================
%  REWRITTEN: Abstract + Introduction
%  Target: JAIR (Journal of Artificial Intelligence Research)
%  Changes from AAAI version:
%   - Abstract starts from application problem, not WMC terminology
%   - Introduction has 7-paragraph arc accessible to ML readers
%   - All technical terms defined on first use
%   - Paper organization paragraph added
% ============================================================

% ────────────────────────────────────────────────────────────
\begin{abstract}
% ────────────────────────────────────────────────────────────
Neural networks that assign scores to concepts governed by a symbolic
ontology face a consistency problem at inference time: a protein may
receive a higher predicted probability of belonging to a specialized
cellular-component term than to the broader term containing it, even
though the ontology mandates that any member of the narrower class must
also be a member of the broader one. Weighted model counting (WMC)
addresses this by treating predictor scores as Bernoulli parameters,
summing the weights of all assignments that satisfy the ontology, and
renormalizing. The result is a principled posterior distribution whose
support is confined to logically valid labelings, computed entirely from
fixed predictor scores without any retraining.

The computational bottleneck is structural width: exact WMC is
exponential in the treewidth of the constraint graph, and large
biomedical ontologies contain thousands of reconvergent regions where
independent messages from different paths double-count shared evidence,
making exact inference intractable for monolithic solvers and path-wise
approximations inaccurate.

We introduce \emph{modular separator WMC}, a framework that
decomposes the constraint graph into a small core and peripheral
modules, eliminates each module into a joint boundary factor (separator),
and contracts a reduced core. Exact separators provably recover the
global WMC posterior. For wide boundaries, we construct nonnegative
tensor-train (TT) separators by querying a boundary-conditioned
oracle directly, project logical support separately, and inject the
factorized TT into reduced-core inference without materializing dense
tables. Recursive core contraction handles ontologies whose reduced
core remains intractable at the first level.

We prove that exact separators preserve the monolithic WMC posterior
(Proposition~\ref{prop:exact-modular}) and that module marginal error
is bounded by boundary posterior error under TV distance
(Proposition~\ref{prop:boundary-error}). Experiments on five complete
biomedical ontologies---Gene Ontology cellular component, molecular
function, and biological process; UBERON anatomy; and ChEBI chemistry,
comprising up to 205,592 terms---show exact marginal recovery on
instances where monolithic references are available, completion of
all-term inference beyond the 300\,GB memory limit of monolithic
junction-tree inference, and systematic reduction in recovered-marginal
error as separator rank increases. On the CAFA Critical Assessment of
Function Annotation benchmark, ontology-conditioned posterior correction
eliminates true-path violations and improves calibration over raw
predictor scores.
\end{abstract}


% ────────────────────────────────────────────────────────────
\section{Introduction}
\label{sec:intro}
% ────────────────────────────────────────────────────────────

% ── ¶1  Application context: the consistency problem ──────
Neuro-symbolic prediction tasks pair a neural scorer with a symbolic
knowledge base that restricts which output configurations are valid.
In hierarchical multi-label classification, for example, a classifier
assigns independent scores to labels arranged in a taxonomy, but the
taxonomy itself demands that any positive prediction at a specific
concept implies positive predictions at all of its ancestors
\citep{silla2011survey}. In protein-function prediction, the Gene
Ontology (GO) encodes is-a relations among hundreds of thousands of
molecular-function, biological-process, and cellular-component terms:
a protein predicted to perform a highly specialized enzymatic activity
must also be predicted to perform the general biochemical processes
containing it \citep{ashburner2000go,geneontology2023}. More broadly,
ontology-based classification---across anatomy \citep{mungall2012uberon},
chemistry \citep{hastings2016chebi}, phenotype, and disease---always faces
the same issue: a vector of local predictor scores is inconsistent with
probability one, because independent scoring ignores the implication
structure encoded in the ontology. We call any complete assignment that
satisfies all ontology constraints a \emph{valid labeling}; raw
predictor scores typically assign nonzero probability to millions of
invalid labelings.

% ── ¶2  WMC as the principled fix ─────────────────────────
\emph{Weighted model counting} (WMC) provides a principled, training-free
mechanism to correct this inconsistency. Given a Boolean formula $T$
(encoding the ontology constraints) and nonnegative literal weights
(derived from predictor scores), WMC sums the weights of all satisfying
assignments of $T$ \citep{chavira2008wmc}. Interpreting the predictor
output $p_i \in (0,1)$ for ontology term $i$ as the Bernoulli parameter
of $y_i \in \{0,1\}$, the product reference distribution
$q_p(y) = \prod_i p_i^{y_i}(1-p_i)^{1-y_i}$ has support over all $2^m$
assignments. Normalizing the restriction of $q_p$ to valid labelings
defines the \emph{WMC posterior}
\[
  p_T(y) = \frac{\mathbf{1}[y \models T]\,q_p(y)}{Z_T(p)},
  \qquad
  Z_T(p) = \mathrm{WMC}(T,p) = \sum_{y \models T} q_p(y).
\]
Every marginal $\mu_i = p_T(y_i=1)$ and every Boolean query uses the
same normalizer $Z_T(p)$, so the corrected scores are globally
consistent by construction. This \emph{inference-time} correction
requires no modification to the predictor architecture, training
procedure, or loss function: it is a plug-in posterior layer applied
to fixed scores at test time.

% ── ¶3  The scalability gap ────────────────────────────────
Despite its conceptual appeal, exact WMC is \#P-hard in general
\citep{valiant1979counting}, and practical solvers based on knowledge
compilation or variable elimination can require time and memory
exponential in the \emph{treewidth} of the constraint graph
$G_T$---the graph whose nodes are propositional variables and whose
edges connect variables that co-occur in a clause
\citep{bodlaender1993treewidth,darwiche2002kcmap,darwiche2011sdd}.
State-of-the-art WMC solvers such as ADDMC, DPMC, and d4
\citep{dudek2020addmc,dudek2020dpmc,lagniez2017improved} achieve
remarkable scalability on random or industrial instances, but biomedical
ontologies expose a different structural regime: the number of variables
is orders of magnitude larger than in benchmark suites, while the
logical structure---is-a implications forming a directed acyclic
graph---is comparatively simple. GO biological process, for example,
encodes 25,950 terms with thousands of reconvergent paths where
multiple specialized terms share a common ancestor through different
routes. On such inputs, the junction-tree backend exhausts a 300\,GB
memory budget before completing even a single ontology.

The cheaper alternative---\emph{path-wise propagation} or loopy belief
propagation (BP) \citep{pearl1988probabilistic}---passes messages along
individual paths through the ontology graph and merges them
independently at each variable. This is computationally efficient but
formally incorrect on reconvergent structures. When two paths from
a common ancestor $d$ arrive at a shared descendant $a$ through
intermediate variables $b$ and $c$, messages along $d \to b \to a$
and $d \to c \to a$ both carry evidence from $d$; treating them as
independent double-counts that evidence and changes the posterior
defined by global WMC. We make this failure mode precise in
Section~\ref{sec:local-reconvergence} with a four-variable example
that shows the path-wise estimate and the exact WMC marginal diverging
by more than a relative error of $10\%$. The scalability challenge is
therefore twofold: avoid the exponential cost of monolithic exact
inference, while retaining the dependence structure required by the
conditioned-WMC posterior.

% ── ¶4  Why existing approaches are insufficient ──────────
Compositional WMC methods \citep{dal2017partitioned,dal2021compositional}
reduce cost by partitioning a formula into components that communicate
through separator variables. These methods were designed for static
formula decomposition and do not address the regime where boundary sizes
routinely reach hundreds or thousands of variables, as observed in GO
biological process (max $|B_j| = 624$) and ChEBI ($|B_j| = 1{,}082$)
in our experiments. Dense separator tables are infeasible at these
widths, and no practical compact representation for the joint boundary
factor has been proposed in the WMC setting.
Low-rank tensor representations have been applied to graphical model
inference \citep{novikov2014mrf,wrigley2017tensorbp} and knowledge
compilation \citep{goral2024semiring,onaka2025ttkc}, but direct
construction of a TT separator from boundary-conditioned oracle queries---
without compiling the full dense table first---is not covered by
existing frameworks.

% ── ¶5  Our approach ──────────────────────────────────────
We introduce \emph{modular separator WMC} to address both requirements
simultaneously. Our framework proceeds in three stages.

\textbf{Boundary-aware decomposition.}
A core-selection algorithm identifies a small set of variables $C$ such
that the remaining variables partition into \emph{modules}---connected
components of $G_T \setminus C$---each with a boundary-constrained
width within a prescribed cap. The algorithm is greedy, adding
high-degree variables to $C$ in batches until all module widths are
manageable, and terminates with a recursive fallback when the reduced
core itself is too wide.

\textbf{Module elimination and separator construction.}
Each module $M_j$ is eliminated by summing out its interior while
holding boundary variables $B_j$ fixed, producing a \emph{separator
factor} $F_j(y_{B_j})$---a single joint function that captures all
dependence the module induces on the rest of the ontology. Crucially,
$F_j$ is defined over the full joint assignment to $B_j$: when
$|B_j| > 1$, it retains correlations among boundary variables that a
factored representation would lose, exactly as required by the
reconvergence analysis. For boundaries of manageable size, $F_j$ is
stored as a dense table. For wide boundaries, we build a
\emph{nonnegative tensor-train separator} by querying the factor at
sampled boundary assignments and fitting a rank-$r$ TT using
alternating nonnegative least squares. Logical support---which boundary
assignments are reachable---is maintained separately and injected as
Boolean constraints, ensuring that no probability mass is placed on
logically unsatisfiable interfaces regardless of approximation error.

\textbf{Core contraction and marginal recovery.}
The original normalizer is recovered by contracting the reduced core with
one separator factor per eliminated module. When every separator is
exact, this equals monolithic WMC (Proposition~\ref{prop:exact-modular}).
A calibrated module junction tree then reuses its messages to recover
every eliminated variable's marginal through backward substitution.
When TT separators are used, the TT cores are injected as local factors
into the reduced-core factor graph, so contraction operates directly on
the factorized representation without materializing the dense table.

% ── ¶6  Contributions ─────────────────────────────────────
\paragraph{Contributions.}
\begin{enumerate}
  \item We introduce modular separator WMC, which decomposes the
  constraint graph into a small core and peripheral modules, eliminates
  each module into a joint separator factor, and recovers every
  ontology-term marginal through backward calibration. The framework
  handles exact and approximate separators through a common interface.

  \item We construct wide-boundary separators directly from
  boundary-conditioned module-WMC queries using a nonnegative
  tensor-train fit, projecting logical support separately to preserve
  constraint correctness independent of approximation rank.

  \item We prove exact modular contraction
  (Proposition~\ref{prop:exact-modular}) and stability of module
  recovery under approximate boundary posteriors
  (Proposition~\ref{prop:boundary-error}).

  \item We validate the framework on five complete biomedical ontologies
  with up to 205,592 terms. On instances where the decomposition is
  non-trivial (recursion depth $> 0$, boundary size $> 0$), exact
  separators recover all reference marginals to machine precision. On
  instances where monolithic inference is infeasible, rank-controlled
  TT separators complete all-term inference within practical time and
  memory budgets, and rank self-consistency confirms accuracy stability.
  On the CAFA protein-function benchmark, posterior correction
  eliminates constraint violations and improves calibration.
\end{enumerate}

% ── ¶7  Paper organization ─────────────────────────────────
\paragraph{Paper organization.}
Section~\ref{sec:related-work} reviews related work in constrained
prediction, exact WMC, compositional and tensor-based model counting,
and ontology-based prediction.
Section~\ref{sec:preliminaries} establishes notation and formally
defines the WMC posterior, constraint graphs, treewidth, and the
core-module-boundary decomposition.
Section~\ref{sec:local-reconvergence} motivates joint separators
with a four-variable reconvergence example.
Section~\ref{sec:method} presents the full modular separator WMC
framework, including the core-selection algorithm, exact and TT
separator construction, core contraction, and marginal recovery.
Section~\ref{sec:theory} states and proves the exactness and stability
guarantees.
Section~\ref{sec:experiments} reports experiments covering correctness
validation on non-trivial decompositions, scalability, separator
approximation quality, comparison with WMC solvers and loopy BP, and
downstream protein-function prediction on GO cellular component and
molecular function.
Section~\ref{sec:discussion} discusses the role of ontology structure
in enabling effective decomposition, limitations of the framework, and
directions for future work.
Section~\ref{sec:conclusion} concludes.
